\documentclass[11pt]{article}

\usepackage[margin=1in]{geometry}
\usepackage{amsmath,amssymb,mathtools,microtype}

\title{Exact Kalman Linear Attention:\\The Parallelization Obstruction}
\author{}
\date{}

\newcommand{\vct}[1]{\boldsymbol{#1}}
\newcommand{\tr}{\mathsf{T}}
\newcommand{\R}{\mathbb{R}}
\DeclareMathOperator{\diag}{diag}
\DeclareMathOperator*{\argmin}{arg\,min}

\begin{document}
\maketitle

\section{Exact Kalman update}

Let $\widetilde{\vct{S}}_t\in\R^{d\times d_v}$ denote the latent memory and
$\vct{S}_t$ its filtered mean, with $\vct{k}_t,\vct{q}_t\in\R^d$ and
$\vct{v}_t\in\R^{d_v}$.  We consider the linear--Gaussian model
\begin{align}
  \widetilde{\vct{S}}_t
  &=\vct{D}_t\widetilde{\vct{S}}_{t-1}+\vct{W}_t,
  &
  \vct{W}_t&\sim\mathcal N_{\mathrm{col}}(\vct{0},\vct{\Omega}_t),
  \\
  \vct{v}_t
  &=\widetilde{\vct{S}}_t^{\tr}\vct{k}_t+\vct{e}_t,
  &
  \vct{e}_t&\sim\mathcal N(\vct{0},r_t\vct{I}),
  \qquad r_t>0.
\end{align}
We assume diagonal transition and process-noise covariance,
\begin{equation}
  \vct{D}_t=\diag(\vct{d}_t),
  \qquad
  \vct{\Omega}_t=\diag(\vct{\omega}_t),
  \qquad \vct{\omega}_t\geq\vct{0}.
\end{equation}
Here $\mathcal N_{\mathrm{col}}(\vct{0},\vct{\Omega}_t)$ means that the
columns of $\vct{W}_t$ are independent $\mathcal N(\vct{0},\vct{\Omega}_t)$
variables; process and observation noises are independent across time.  We
take $\vct{P}_0\succ\vct{0}$ and assume each predictive covariance below is
positive definite.
The posterior covariance $\vct{P}_t\in\R^{d\times d}$ is nevertheless dense:
conditioning on a general key couples its coordinates.

Given $(\vct{S}_{t-1},\vct{P}_{t-1})$, the exact Kalman filter is
\begin{align}
  \widehat{\vct{S}}_t
  &=\vct{D}_t\vct{S}_{t-1},
  &
  \widehat{\vct{P}}_t
  &=\vct{D}_t\vct{P}_{t-1}\vct{D}_t^{\tr}+\vct{\Omega}_t,
  \\
  \vct{\kappa}_t
  &=\frac{\widehat{\vct{P}}_t\vct{k}_t}
  {r_t+\vct{k}_t^{\tr}\widehat{\vct{P}}_t\vct{k}_t},
  &
  \vct{S}_t
  &=\widehat{\vct{S}}_t+
  \vct{\kappa}_t
  \bigl(\vct{v}_t-\widehat{\vct{S}}_t^{\tr}\vct{k}_t\bigr)^{\tr},
  \\
  \vct{P}_t
  &=\widehat{\vct{P}}_t-
  \frac{(\widehat{\vct{P}}_t\vct{k}_t)
  (\widehat{\vct{P}}_t\vct{k}_t)^{\tr}}
  {r_t+\vct{k}_t^{\tr}\widehat{\vct{P}}_t\vct{k}_t}.
\end{align}
The corresponding linear-attention read is
$\vct{o}_t=\vct{S}_t^{\tr}\vct{q}_t$.

\section{Information form}

Define the posterior and predictive precision matrices
\begin{equation}
  \vct{C}_t=\vct{P}_t^{-1},
  \qquad
  \widehat{\vct{C}}_t=\widehat{\vct{P}}_t^{-1}.
\end{equation}
The Woodbury identity converts the rank-one covariance correction into the
additive information update
\begin{equation}
  \boxed{
  \vct{C}_t
  =\widehat{\vct{C}}_t+\frac{\vct{k}_t\vct{k}_t^{\tr}}{r_t}}
  =
  \left(
    \vct{D}_t\vct{C}_{t-1}^{-1}\vct{D}_t^{\tr}
    +\vct{\Omega}_t
  \right)^{-1}
  +\frac{\vct{k}_t\vct{k}_t^{\tr}}{r_t}.
  \label{eq:exact-information}
\end{equation}
Equivalently, with the whitened key
$\bar{\vct{k}}_t=\vct{k}_t/\sqrt{r_t}$, the measurement contributes
$\bar{\vct{k}}_t\bar{\vct{k}}_t^{\tr}$.

Under the static, unit-noise specialization
$\vct{D}_t=\vct{I}$, $\vct{\Omega}_t=\vct{0}$, and $r_t=1$, this reduces to
\begin{equation}
  \vct{C}_t=\vct{C}_{t-1}+\vct{k}_t\vct{k}_t^{\tr}
  \label{eq:static-information}
\end{equation}
If $\vct{\Omega}_t=\vct{0}$ and $\vct{D}_t$ is invertible, the more general
update is
\begin{equation}
  \vct{C}_t
  =\vct{D}_t^{-\tr}\vct{C}_{t-1}\vct{D}_t^{-1}
  +\frac{\vct{k}_t\vct{k}_t^{\tr}}{r_t}.
\end{equation}
With generic nonzero process noise, even diagonal $\vct{D}_t$ and
$\vct{\Omega}_t$ leave the prediction in Equation~\eqref{eq:exact-information}
as a nonlinear dense matrix map.

\section{Why parallelization remains difficult}

The static case makes the distinction particularly clear.  Define the
information target $\vct{B}_t=\vct{C}_t\vct{S}_t$.  Then
\begin{equation}
  \vct{C}_t
  =\vct{C}_{t-1}+\frac{\vct{k}_t\vct{k}_t^{\tr}}{r_t},
  \qquad
  \vct{B}_t
  =\vct{B}_{t-1}+\frac{\vct{k}_t\vct{v}_t^{\tr}}{r_t}.
\end{equation}
Both are associative prefix sums.  The difficult operation is the causal
readout
\begin{equation}
  \vct{S}_t=\vct{C}_t^{-1}\vct{B}_t,
  \qquad
  \vct{\kappa}_t=\frac{\vct{C}_t^{-1}\vct{k}_t}{r_t},
\end{equation}
which requires a different dense solve at every prefix.  Factoring only the
final precision would expose earlier outputs to future keys.  Sequential
rank-one inverse updates have depth $T$.  Factoring every prefix independently
exposes concurrency, but repeats a dense factorization $T$ times rather than
sharing work through a useful prefix operator.

For the dynamic model, the problem is harder still: process noise makes the
precision prediction itself non-additive.  Thus the obstacle is not forming
$\vct{k}_t\vct{k}_t^{\tr}$, but applying a prefix-dependent dense inverse at
every token while preserving causality.  With invertible transitions the
Riccati map admits a formal associative lifting to $2d\times2d$
matrix-fractional maps, but each composition is a dense matrix operation and
the products become numerically stiff under strong contraction.  The
algorithmic question is therefore not whether a formal parallel prefix exists,
but whether its composition can be made inexpensive enough to replace the
tokenwise chain.

\section{Diagonal information and online variational inference}

Suppose instead that the retained precision is diagonal,
\begin{equation}
  \vct{C}_{t-1}=\diag(\vct{c}_{t-1}),
  \qquad
  \vct{C}_{t-1}^{-1}
  =\diag(\vct{c}_{t-1}^{-1}).
\end{equation}
Because $\vct{D}_t$ and $\vct{\Omega}_t$ are also diagonal, the predictive
precision remains diagonal:
\begin{equation}
  \widehat{\vct{C}}_t=\diag(\widehat{\vct{c}}_t),
  \qquad
  \widehat{\vct{c}}_t
  =\frac{\vct{c}_{t-1}}
  {\vct{d}_t^2+\vct{\omega}_t\odot\vct{c}_{t-1}},
  \label{eq:diagonal-information-prediction}
\end{equation}
where all vector operations are elementwise.  The exact gain under this
predictive state is
\begin{equation}
  \vct{\kappa}_t
  =\frac{\widehat{\vct{c}}_t^{-1}\odot\vct{k}_t}
  {r_t+\sum_i k_{t,i}^2/\widehat c_{t,i}},
  \label{eq:diagonal-information-gain}
\end{equation}
so the uncertainty recurrence decomposes into independent scalar maps.  All
of their prefixes can be generated by an associative scan, after which the
gain is evaluated tokenwise without a dense prefix solve.

The exact measurement update does not preserve this structure.  By
Equation~\eqref{eq:exact-information}, its precision is
\begin{equation}
  \vct{C}_t^\star
  =\diag(\widehat{\vct{c}}_t)
   +\frac{\vct{k}_t\vct{k}_t^{\tr}}{r_t},
  \label{eq:exact-dense-posterior-information}
\end{equation}
which is generally dense.  We therefore project the exact one-step posterior
$\pi_t$ back to the diagonal family after every token:
\begin{equation}
  q_t
  =\argmin_{q=\mathcal N_{\mathrm{col}}
    (\vct{M},\diag(\vct{c}^{-1}))}
    \operatorname{KL}(q\,\|\,\pi_t).
  \label{eq:online-diagonal-vi}
\end{equation}
The reverse-KL optimum preserves the exact posterior mean,
$\vct{M}=\vct{S}_t$, and retains the diagonal of the exact posterior
precision:
\begin{equation}
  \boxed{
  \vct{c}_t
  =\widehat{\vct{c}}_t
   +\frac{\vct{k}_t\odot\vct{k}_t}{r_t}.}
  \label{eq:diagonal-vi-information}
\end{equation}
Feeding $\diag(\vct{c}_t^{-1})$ into the next prediction closes the online
recursion.  With $u_{t,i}=k_{t,i}^2/r_t$, each channel obeys
\begin{equation}
  c_{t,i}
  =\frac{c_{t-1,i}}
  {d_{t,i}^2+\omega_{t,i}c_{t-1,i}}+u_{t,i}.
  \label{eq:diagonal-mobius}
\end{equation}
This is a scalar M\"obius map, whose token-local $2\times2$ representations
compose associatively.  Online variational inference therefore replaces each
prefix-dependent dense inverse with $d$ independent scalar scans while
retaining the exact Kalman mean update under the diagonal predictive prior.

\section{Oracle-free gain approximation}

The gain need not be computed by first forming a covariance inverse.  Combining
the information update with the Kalman gain gives the equivalent linear system
\begin{equation}
  \boxed{
  \vct{C}_t\vct{\kappa}_t=\frac{\vct{k}_t}{r_t},
  \qquad
  \vct{C}_t
  =\widehat{\vct{C}}_t
   +\frac{\vct{k}_t\vct{k}_t^{\tr}}{r_t}.}
  \label{eq:gain-linear-system}
\end{equation}
This suggests approximating the solution of a known positive-definite system,
rather than defining an approximation in terms of the unknown exact gain.

\subsection{Diagonal-Kalman-compatible elementwise gains}

Consider
\begin{equation}
  \vct{\kappa}_t(\vct{b})=\vct{b}\odot\vct{k}_t.
  \label{eq:elementwise-gain-family}
\end{equation}
Let $\mathcal I_t=\{i:k_{t,i}\neq0\}$ and set $b_i=0$ outside
$\mathcal I_t$.  For a positive diagonal predictive precision,
Equation~\eqref{eq:diagonal-information-gain} gives
\begin{equation}
  b_i
  =\frac{\widehat c_{t,i}^{-1}}{r_t+\zeta_t},
  \qquad
  \zeta_t
  =\sum_{j\in\mathcal I_t}\frac{k_{t,j}^2}{\widehat c_{t,j}}.
  \label{eq:kalman-channel-coefficient}
\end{equation}
Thus $b_i\geq0$, and the effective write along the current key is
\begin{equation}
  \beta_t^{\mathrm{eff}}
  :=\vct{k}_t^{\tr}\vct{\kappa}_t
  =\sum_{i\in\mathcal I_t}b_i k_{t,i}^2
  =\frac{\zeta_t}{r_t+\zeta_t}
  \in[0,1).
  \label{eq:effective-same-key-write}
\end{equation}
If $\widehat{\vct{y}}_t=\widehat{\vct{S}}_t^{\tr}\vct{k}_t$ is the value
recalled before the update, then
\begin{align}
  \vct{S}_t^{\tr}\vct{k}_t
  &=\widehat{\vct{y}}_t
    +\beta_t^{\mathrm{eff}}
      (\vct{v}_t-\widehat{\vct{y}}_t),
  \\
  \vct{v}_t-\vct{S}_t^{\tr}\vct{k}_t
  &=(1-\beta_t^{\mathrm{eff}})
    (\vct{v}_t-\widehat{\vct{y}}_t).
  \label{eq:same-key-residual-contraction}
\end{align}
The first constraint below prevents a channel from writing opposite to its key
coordinate, while the second bounds the total same-key write:
\begin{equation}
  \mathcal B_t
  :=\left\{\vct{b}:
  b_i\geq0\ (i\in\mathcal I_t),\quad
  \sum_{i\in\mathcal I_t}b_i k_{t,i}^2\leq1\right\}.
  \label{eq:diagonal-gain-constraints}
\end{equation}
The finite positive-precision family satisfies strict inequalities.  The weak
inequalities include its closure, and every interior point corresponds to the
positive diagonal predictive precision
\begin{equation}
  \widehat c_i
  =\frac{1-\sum_{j\in\mathcal I_t}b_jk_{t,j}^2}{r_tb_i},
  \qquad i\in\mathcal I_t.
  \label{eq:gain-to-diagonal-information}
\end{equation}

There is an important limitation: if every key coordinate is nonzero, the map
$\vct{b}\mapsto\vct{b}\odot\vct{k}_t$ is a bijection.  An unrestricted
channelwise $\vct{b}$ can therefore represent any gain and is no easier to find
than the original dense solve.  Computational savings require either a
lower-dimensional structure on $\vct{b}$ or a fixed-work approximate solver.

\subsection{A computable variational problem}

Let $\vct{K}_t=\diag(\vct{k}_t)$ and
$\vct{s}_t=\vct{k}_t\odot\vct{k}_t$.  Restricting the quadratic variational
form of Equation~\eqref{eq:gain-linear-system} to
$\vct{\kappa}_t=\vct{K}_t\vct{b}$ gives
\begin{equation}
  \boxed{
  \vct{b}_t^\dagger
  \in\argmin_{\vct{b}\in\mathcal B_t}
  \Phi_t(\vct{b}),
  \qquad
  \Phi_t(\vct{b})
  =\frac{r_t}{2}\vct{b}^{\tr}\vct{K}_t\vct{C}_t\vct{K}_t\vct{b}
   -\vct{s}_t^{\tr}\vct{b}.}
  \label{eq:computable-gain-objective}
\end{equation}
Unlike an oracle projection, this objective contains only the available
precision, key, and observation noise.  Its gradient is
\begin{equation}
  \nabla\Phi_t(\vct{b})
  =r_t\vct{k}_t\odot
    \bigl[\vct{C}_t(\vct{b}\odot\vct{k}_t)\bigr]
   -\vct{s}_t.
  \label{eq:computable-gain-gradient}
\end{equation}
Without the constraints and with a full-support key, its stationarity equation
is exactly Equation~\eqref{eq:gain-linear-system}; with the constraints it is
the $\vct{C}_t$-energy projection of that system onto the capped
diagonal-compatible gain set.  This is a computable energy surrogate, not the
constrained posterior-MSE projection: the two objectives share the exact
unconstrained minimizer but use different projection metrics.  Evaluating
Equation~\eqref{eq:computable-gain-objective} never requires
$\widehat{\vct{C}}_t^{-1}\vct{k}_t$.

\paragraph{A low-dimensional coefficient family.}
One explicit structural restriction is
\begin{equation}
  \vct{b}=\vct{R}_t\vct{\theta},
  \qquad \vct{R}_t\in\R^{d\times m},
  \qquad m\ll d,
  \label{eq:tied-coefficient-family}
\end{equation}
where columns of $\vct{R}_t$ can, for example, tie channels into fixed groups.
The computable problem becomes the $m$-dimensional quadratic program
\begin{equation}
  \vct{\theta}_t^\dagger
  \in\argmin_{\substack{\vct{R}_t\vct{\theta}\geq\vct{0}\\
               \vct{s}_t^{\tr}\vct{R}_t\vct{\theta}\leq1}}
  \left\{
  \frac12\vct{\theta}^{\tr}\vct{H}_t\vct{\theta}
  -\vct{h}_t^{\tr}\vct{\theta}\right\},
  \label{eq:grouped-gain-problem}
\end{equation}
where
\begin{equation}
  \vct{H}_t
  =r_t\vct{R}_t^{\tr}\vct{K}_t\vct{C}_t\vct{K}_t\vct{R}_t,
  \qquad
  \vct{h}_t=\vct{R}_t^{\tr}\vct{s}_t.
\end{equation}
This needs $m$ precision--vector products followed by a small local solve.  For
$\vct{k}_t\neq\vct{0}$, taking $m=1$ and $\vct{R}_t=\vct{1}$ gives
$\vct{\kappa}_t=\beta_t\vct{k}_t$ with the closed form
\begin{equation}
  \boxed{
  \beta_t
  =\frac{\|\vct{k}_t\|_2^2}
  {r_t\vct{k}_t^{\tr}\vct{C}_t\vct{k}_t},
  \qquad
  \vct{k}_t^{\tr}\vct{\kappa}_t
  =\frac{\|\vct{k}_t\|_2^4}
  {r_t\vct{k}_t^{\tr}\widehat{\vct{C}}_t\vct{k}_t
   +\|\vct{k}_t\|_2^4}<1.}
  \label{eq:scalar-galerkin-gain}
\end{equation}
The first expression uses the posterior precision $\vct{C}_t$; the second uses
$\vct{C}_t=\widehat{\vct{C}}_t+\vct{k}_t\vct{k}_t^{\tr}/r_t$ to expose the
same-key bound.  A nested sequence of column spaces for $\vct{R}_t$
interpolates between this scalar gain and the full channelwise problem.
Setting $m=d$ and $\vct{R}_t=\vct{I}$ removes the structural saving and
recovers the full constrained channelwise problem; it equals the dense Kalman
gain only when that gain belongs to $\mathcal B_t$.

\paragraph{A fixed-work channelwise solver.}
Alternatively, retain one coefficient per channel but run exactly $L$
projected-gradient steps,
\begin{equation}
  \vct{b}_t^{(\ell+1)}
  =\Pi_{\mathcal B_t}\!\left(
    \vct{b}_t^{(\ell)}
    -\eta_{t,\ell}
      \left\{
      r_t\vct{k}_t\odot
       [\vct{C}_t(\vct{b}_t^{(\ell)}\odot\vct{k}_t)]
      -\vct{s}_t
      \right\}
  \right),
  \qquad \ell=0,\ldots,L-1.
  \label{eq:projected-gain-iteration}
\end{equation}
The projection onto $\mathcal B_t$ is a nonnegative weighted-simplex
projection, implemented by a scalar bisection or sorting.  A sufficient fixed
step-size condition is
\begin{equation}
  0<\eta_{t,\ell}\leq
  \lambda_{\max}^{-1}\!\left(
    r_t\vct{K}_t\vct{C}_t\vct{K}_t
  \right),
  \label{eq:projected-gain-step-size}
\end{equation}
with the matrices restricted to the support of $\vct{k}_t$; a monotone line
search is another option.  Equation~\eqref{eq:kalman-channel-coefficient}
provides a natural diagonal-variational initialization for $\vct{b}_t^{(0)}$.
The integer $L$ is now an explicit speed--accuracy knob: every iteration uses
one dense matrix--vector product, reductions, and elementwise operations, but
no inverse and no oracle gain.

\subsection{Parallel scope}

In the static model, all posterior precisions are available from the
associative prefix
\begin{equation}
  \vct{C}_t=\vct{C}_0+
  \sum_{\tau\leq t}\frac{\vct{k}_\tau\vct{k}_\tau^{\tr}}{r_\tau}.
\end{equation}
At projected-gradient iteration $\ell$, the required product for every token
can equivalently be written
\begin{equation}
  \vct{C}_t\vct{u}_t^{(\ell)}
  =\vct{C}_0\vct{u}_t^{(\ell)}
   +\sum_{\tau\leq t}\frac{\vct{k}_\tau
      (\vct{k}_\tau^{\tr}\vct{u}_t^{(\ell)})}{r_\tau},
  \qquad
  \vct{u}_t^{(\ell)}=\vct{b}_t^{(\ell)}\odot\vct{k}_t.
  \label{eq:prefix-precision-matvec}
\end{equation}
Thus all positions can execute each solver stage concurrently.  Materializing
all dense prefix matrices once gives sequence depth $O(\log T+L)$, after which
the $L$ matrix--vector steps are token-local.  Alternatively,
Equation~\eqref{eq:prefix-precision-matvec} is a causal linear-attention
contraction with query $\vct{u}_t^{(\ell)}$; recomputing it at every stage gives
$O(L\log T)$ sequence depth without materializing every prefix matrix.  In
either schedule, the construction replaces each dense inverse by a fixed
number of parallel precision--vector products.  It does not remove dense
arithmetic, and with nonzero process noise constructing the exact dense
$\vct{C}_t$ prefixes still has the nonlinear Riccati obstruction in
Equation~\eqref{eq:exact-information}.

\end{document}
